\documentclass{situationist-preprint}

\title{The Scope of Context-Sensitivity}
\msid{LNG-19-a}
\msfamily{LNG-19}
\msdiscipline{Philosophy of Language}
\mstype{Principal manuscript}
\mskeywords{context-sensitivity; indexicals; demonstratives; gradable adjectives; definite descriptions; quantifier domain restriction; minimalism; contextualism}

\begin{document}
\maketitle

\begin{abstract}
Almost every expression has been claimed to be context-sensitive: quantifiers, definite descriptions, colour words, \enquote{know}, \enquote{ready}, \enquote{rain}. If all these claims were right, the semantic content of a sentence would be fixed only by the whole speech situation, and the distinction between what a sentence says and what a speaker conveys with it would lose its point. We argue that context-sensitivity extends only to indexicals, demonstratives and gradable adjectives, together with expressions whose lexical analysis contains one of these. The criterion is semantic incompleteness: an expression is context-sensitive when its standing meaning determines no content until context supplies a value for a parameter that the meaning itself introduces. Indexicals, demonstratives and the standard of a gradable adjective meet the criterion. Definite descriptions do not: \enquote{the table} has a complete Russellian content in every context, and the restriction that makes an incomplete description apt is a feature of what the speaker communicates, not of what the description means. The same holds of quantifier domains generally. We answer the binding argument, the argument from truth-value intuitions, and the objection that \enquote{know} would require a new category, and we show how the account treats colour words and predicates of personal taste.
\end{abstract}

\section{The problem}

That some expressions are context-sensitive is not in dispute. \enquote{I}, \enquote{here}, \enquote{now} and \enquote{today} have a fixed linguistic meaning, which \citet{Kaplan1989} called their character, and a content that varies with the context of utterance; \enquote{this} and \enquote{that} have a content fixed by a demonstration or by the speaker's directing intention. The dispute concerns how far beyond these the phenomenon extends.

Radical contextualists hold that it extends everywhere. \citet{Searle1980} argued that literal meaning determines truth conditions only against a non-semantic background, and \citet{Travis1997} constructed cases in which one and the same sentence, \enquote{The leaves are green}, uttered of leaves painted green, says something true in one context and something false in another. \citet{Recanati2004} holds that what is said is the output of pragmatic processes of enrichment and modulation that operate on lexical meanings before any truth-evaluable content exists. At the other extreme, \citet{CappelenLepore2005} defend a minimal semantics on which only a short list of expressions, the \enquote{basic set}, is context-sensitive, and the variability on which contextualists rely belongs to the many propositions a speaker asserts, not to the one a sentence semantically expresses. Between the two, \citet{Stanley2000} argues that all genuine effects of context on truth conditions are traceable to covert variables in logical form, and \citet{StanleySzabo2000} locate the domain restriction of quantifiers in such a variable attached to the nominal.

Definite descriptions sit at the centre of this dispute. On the theory of \citet{Russell1905}, which LNG-03 defends, \enquote{the $F$ is $G$} says that there is exactly one $F$ and it is $G$. Since \enquote{The table is covered with books} is usually uttered where there are many tables, the theory seems to require that the description be completed by context, and so to make descriptions context-sensitive after all \citep{Strawson1950,Wettstein1981}. We argue that it does not. Context-sensitivity extends only to indexicals, demonstratives and gradable adjectives; definite descriptions are not context-sensitive, and the restriction that makes an incomplete description apt is pragmatic.

\section{Framework}

\begin{definition}[Standing meaning and content]\label{d:meaning}
The \emph{standing meaning} of an expression is what a competent speaker knows in knowing the expression, independently of any occasion of use. Its \emph{content} on an occasion is its contribution to the proposition the sentence semantically expresses on that occasion.
\end{definition}

\begin{definition}[Semantic incompleteness]\label{d:incomplete}
An expression $e$ is \emph{context-sensitive} if its standing meaning introduces a parameter such that, unless the context supplies a value for it, no sentence in which $e$ occurs free determines a content. An expression whose standing meaning determines a content without any such value is \emph{context-invariant}, whatever else context may do with it.
\end{definition}

\begin{definition}[The three categories]\label{d:categories}
An \emph{indexical} is an expression whose standing meaning is a rule mapping a feature of the context of utterance (agent, time, place, world) to a content. A \emph{demonstrative} is an expression whose standing meaning requires a demonstration or directing intention to fix its content. A \emph{gradable adjective} is an adjective that denotes a relation between objects and degrees on a scale, and whose positive form requires a standard of comparison on that scale \citep{KennedyMcNally2005,Kennedy2007}.
\end{definition}

The tests for gradability are grammatical: gradable adjectives take comparatives (\enquote{taller}), degree modifiers (\enquote{very tall}, \enquote{tall enough}) and measure phrases (\enquote{two metres tall}). Relative gradable adjectives such as \enquote{tall} take their standard from context; absolute ones such as \enquote{flat} and \enquote{full} have a standard fixed at an end point of the scale, and show context-dependence chiefly in what counts as reaching it \citep{Kennedy2007}.

\begin{principle}[Inheritance]\label{p:inherit}
An expression whose lexical analysis contains an indexical, demonstrative or gradable component is context-sensitive in virtue of that component, and in no other way.
\end{principle}

Principle~\ref{p:inherit} lets an expression of no listed category be context-sensitive without enlarging the list. \enquote{Local} contains an indexical argument place, and so does \enquote{yesterday}. Tense is indexical in the same sense: its standing meaning relates the time of the described event to the time of utterance.

\begin{principle}[Pragmatic enrichment]\label{p:enrich}
A speaker who utters a sentence with content $p$ may communicate a proposition $q$ distinct from $p$, which the hearer recovers by reasoning about the speaker's purposes. Such a $q$ belongs to what the speaker means, not to what the sentence says, even when it is what the speaker primarily intends to convey.
\end{principle}

Principle~\ref{p:enrich} is common ground: \citet{Bach1994} calls such contents implicitures, \citet{Recanati2004} places them in what is said, and \citet{CappelenLepore2005} count them among the propositions asserted. The disagreement is whether they belong to semantic content. On our view they do only when a parameter introduced by the standing meaning requires them, as Definition~\ref{d:incomplete} states. The derivation of such contents proceeds by the same Gricean reasoning that LNG-11 shows to underlie conversational implicature, though enrichments are not themselves implicatures.

\section{The argument}

\begin{lemma}\label{l:three}
Indexicals, demonstratives and gradable adjectives are context-sensitive.
\end{lemma}

\begin{proof}
For an indexical, the standing meaning is a rule whose argument is a feature of the context (Definition~\ref{d:categories}); \enquote{I am hungry} determines no proposition until an agent is given. For a demonstrative, the standing meaning requires a demonstration or intention, and \enquote{that is red} determines no proposition until one is supplied. For a relative gradable adjective, the positive form denotes the property of exceeding a standard on the scale; the standing meaning fixes the scale but not the standard, so \enquote{Ana is tall} determines no proposition until a standard is supplied. Each meets Definition~\ref{d:incomplete}.
\end{proof}

In the gradable case the standard is fixed by the comparison class and the interests in play \citep{Kennedy2007}; that its location is vague does not affect the point, since LNG-12 shows that vague predicates have sharp boundaries we cannot know.

\begin{lemma}\label{l:descr}
Definite descriptions are context-invariant.
\end{lemma}

\begin{proof}
On the Russellian analysis defended in LNG-03, \enquote{the $F$ is $G$} expresses the proposition that exactly one thing is $F$ and every $F$ is $G$. The standing meaning of \enquote{the} is a relation between the properties expressed by $F$ and $G$, and introduces no further parameter. Given the contents of $F$ and $G$, the sentence determines a proposition without any value from context. By Definition~\ref{d:incomplete}, the description is context-invariant. If $F$ itself contains an indexical, as in \enquote{the man I met yesterday}, the description has a content only in context, but by Principle~\ref{p:inherit} it is context-sensitive in virtue of \enquote{I} and \enquote{yesterday}, not in virtue of being a description.
\end{proof}

The incompleteness of the description in \enquote{The table is covered with books} does not refute the lemma. Taken at its content, the sentence says that there is exactly one table in the world, and that it is covered with books. That proposition is false; it is a complete proposition nonetheless. What the speaker communicates is something else: that the unique table in a contextually salient domain, the room or the conversation, is covered with books. By Principle~\ref{p:enrich}, that proposition belongs to what the speaker means. This is the line of LNG-03, on which Strawson's objection concerns use and not content: the uses Strawson described are real, but they do not concern what the description means \citep{Kripke1977,Neale1990}.

\begin{lemma}\label{l:quant}
Quantifier domain restriction is pragmatic. The domain of a quantifier phrase is not a parameter introduced by the standing meaning of the determiner or of the noun.
\end{lemma}

\begin{proof}
Suppose, for reductio, that the standing meaning of \enquote{every}, or of the noun it combines with, introduced a domain parameter in the sense of Definition~\ref{d:incomplete}. Then \enquote{Every bottle is empty} would determine no proposition unless context supplied a domain. But it determines one: that every bottle there is is empty. Speakers can assert exactly that without contextual supplementation; the unrestricted reading is what results when nothing is added. Moreover, when a restricted reading is intended, there is often no determinate restriction: a speaker who says \enquote{Every bottle is empty} after a party need not have in mind any one of the bottles in this kitchen, the bottles we bought, or the bottles opened tonight. A parameter introduced by standing meaning must receive a value for the sentence to have content; communicated contents may be indeterminate in this way. So the restriction is not a semantic parameter, and by Principle~\ref{p:enrich} it belongs to what the speaker communicates.
\end{proof}

\citet{Wettstein1981} used the same indeterminacy against the Russellian analysis: no single completion of \enquote{the table} can be identified as what the speaker said, so the description must be referential. The better conclusion is that no completion is what was said, because nothing was added to what was said; speaker's reference, as \citet{Kripke1977} showed, arises equally in a language stipulated to be Russellian.

\begin{theorem}\label{t:scope}
Context-sensitivity extends only to indexicals, demonstratives and gradable adjectives, and to expressions that contain one of them in their lexical analysis. Definite descriptions are not context-sensitive.
\end{theorem}

\begin{proof}
By Lemma~\ref{l:three}, the three categories are context-sensitive, and by Principle~\ref{p:inherit} so are expressions that contain them. The candidates for further context-sensitivity fall into three groups. Quantifier phrases and definite descriptions, which are quantifier phrases on the Russellian analysis, are context-invariant by Lemmas~\ref{l:descr} and~\ref{l:quant}. Expressions with an unfilled argument place, such as \enquote{ready}, \enquote{enough} or \enquote{rain}, have a complete content in which the place is existentially closed: \enquote{Ana is ready} says that she is ready for something, and \enquote{It is raining} says that it is raining at some place at the time of utterance, where the time is supplied by tense, an indexical. The specific completions speakers convey are, as in the quantifier case, often indeterminate and are recovered by Principle~\ref{p:enrich}. The remaining candidates, colour words, predicates of personal taste and \enquote{know}, contain a gradable component, as Section~\ref{s:cases} shows, and are context-sensitive by Principle~\ref{p:inherit}. No candidate requires a fourth category.
\end{proof}

\section{Cases}\label{s:cases}

\emph{Colour words.} \citet{Travis1997}'s case of the painted leaves shows that \enquote{The leaves are green} may be true in one context and false in another. Colour adjectives are gradable: they take comparatives and degree modifiers (\enquote{greener}, \enquote{very green}, \enquote{green enough}). \citet{KennedyMcNally2010} trace their variability to the standard on the scale and to the way the colour of an object's parts is aggregated into the colour of the object, both features of the gradable adjective. The leaves case confirms the account: the variability is real, and located in a gradable adjective.

\emph{Predicates of personal taste.} \enquote{Tasty} and \enquote{fun} are gradable (\enquote{tastier}, \enquote{very fun}), and the standard of the positive form, as well as the experiencer relative to whom the scale is set, is supplied by context. \citet{Lasersohn2005} argued that such predicates are relative to a judge rather than context-sensitive in the ordinary way, because of faultless disagreement. Whatever the right account of the disagreement, the variation lies in a gradable predicate.

\emph{\enquote{Know}.} EPI-13 argues that \enquote{know} is context-sensitive. Its account is an instance of Principle~\ref{p:inherit}, and introduces no new category. On the analysis of EPI-01, knowledge is safe, justified true belief, and \enquote{safe} is a gradable adjective: it takes comparatives and degree modifiers (\enquote{safer}, \enquote{safe enough}), and what counts as safe depends on a contextual standard on a scale, which in the epistemic case is the ordering of worlds by their closeness to the subject's world. \enquote{Know} inherits its context-sensitivity from the gradable notion in its analysis. \citet{Stanley2004} objected that \enquote{know} is not gradable, since it does not take comparatives or degree modifiers. The observation is correct about the verb, and EPI-13 accepts it; Principle~\ref{p:inherit} does not require the containing expression to be gradable, only that its analysis contain a gradable component.

\section{Objections and replies}

\begin{objection}[Binding]\label{o:bind}
In \enquote{In every room, every bottle is empty}, the domain of \enquote{every bottle} varies with the rooms: the sentence means that in each room, every bottle in that room is empty. A domain that can be bound by a higher quantifier must be represented by a variable in logical form. So quantifier domain restriction is semantic \citep{StanleySzabo2000,Stanley2000}.
\end{objection}

\begin{reply}
The argument overgenerates. In \enquote{Everywhere I go, two plus two is four}, the binding test would posit a location variable in the logical form of an arithmetical sentence \citep{CappelenLepore2002}. Bound readings can be produced when the higher quantifier introduces a variable that the communicated proposition exploits, as \citet{Recanati2004} argues, and nothing in the data requires the variable to be present in the standing meaning of the lower phrase. Moreover, the bound reading of the example is available because the prepositional phrase supplies an explicit restriction, which our account treats as part of content: \enquote{in every room} quantifies over rooms and restricts the embedded clause to each room in turn. What remains pragmatic is restriction without any such linguistic material.
\end{reply}

\begin{objection}[Truth-value judgments]\label{o:intuit}
Speakers judge an utterance of \enquote{Every bottle is empty} after a party to be true when the bottles in the kitchen are empty. If the sentence said that every bottle in the world is empty, they would be judging a false proposition true. The account imputes systematic error to competent speakers.
\end{objection}

\begin{reply}
Speakers' judgments of truth on an occasion track what the speaker said in the broad sense, the proposition she primarily intended to convey, and this is what they ought to track when assessing a speaker. The account imputes no error, since it agrees that the speaker said something true in that sense. What it denies is that the judgment reports the semantic content of the sentence. When the unrestricted reading is made salient (\enquote{Strictly speaking, not every bottle is empty}), speakers recognize a correction of the words, not a change of topic \citep{CappelenLepore2005}.
\end{reply}

\begin{objection}[Ad hoc gradability]\label{o:adhoc}
The account deals with hard cases by finding a gradable adjective in them. With enough ingenuity, any context-sensitivity could be redescribed as the context-sensitivity of some gradable notion, and the thesis becomes unfalsifiable.
\end{objection}

\begin{reply}
Gradability has independent grammatical tests, and the account is falsified by any expression that varies in content with context, is neither indexical nor demonstrative, and neither passes these tests nor contains, on independent grounds, an expression that does. Colour words and taste predicates pass the tests directly. For \enquote{know}, the gradable component is located by an analysis of knowledge argued for on independent grounds in EPI-01, not chosen to save the thesis.
\end{reply}

\section{What the result does not establish}

The argument does not settle how the standards of gradable adjectives are fixed. It does not treat epistemic modals; a treatment consistent with the theorem would have to locate their dependence on a body of information in an indexical or gradable component, and we do not claim to have shown that this can be done. And it does not deny that communicated content varies with context far more than semantic content does: the thesis concerns only the latter. Incomplete descriptions are treated further in LNG-19-b.

\section{Conclusion}

Context-sensitivity, in the sense of a parameter introduced by standing meaning and filled by context, belongs to indexicals, demonstratives and gradable adjectives, and to expressions whose analysis contains them. Quantifier phrases, and with them definite descriptions, have complete contents in every context; the domain restrictions and completions that make their ordinary uses apt are features of what speakers communicate.

\begin{thebibliography}{99}\small
\bibitem[Bach(1994)]{Bach1994} Bach, K. (1994). Conversational impliciture. \emph{Mind \& Language}, 9(2), 124--162.
\bibitem[Cappelen and Lepore(2002)]{CappelenLepore2002} Cappelen, H. and Lepore, E. (2002). Indexicality, binding, anaphora and a priori truth. \emph{Analysis}, 62(4), 271--281.
\bibitem[Cappelen and Lepore(2005)]{CappelenLepore2005} Cappelen, H. and Lepore, E. (2005). \emph{Insensitive Semantics: A Defense of Semantic Minimalism and Speech Act Pluralism}. Blackwell.
\bibitem[Kaplan(1989)]{Kaplan1989} Kaplan, D. (1989). Demonstratives. In J. Almog, J. Perry and H. Wettstein (eds), \emph{Themes from Kaplan}, 481--563. Oxford University Press.
\bibitem[Kennedy(2007)]{Kennedy2007} Kennedy, C. (2007). Vagueness and grammar: The semantics of relative and absolute gradable adjectives. \emph{Linguistics and Philosophy}, 30(1), 1--45.
\bibitem[Kennedy and McNally(2005)]{KennedyMcNally2005} Kennedy, C. and McNally, L. (2005). Scale structure, degree modification, and the semantics of gradable predicates. \emph{Language}, 81(2), 345--381.
\bibitem[Kennedy and McNally(2010)]{KennedyMcNally2010} Kennedy, C. and McNally, L. (2010). Color, context, and compositionality. \emph{Synthese}, 174(1), 79--98.
\bibitem[Kripke(1977)]{Kripke1977} Kripke, S. (1977). Speaker's reference and semantic reference. \emph{Midwest Studies in Philosophy}, 2, 255--276.
\bibitem[Lasersohn(2005)]{Lasersohn2005} Lasersohn, P. (2005). Context dependence, disagreement, and predicates of personal taste. \emph{Linguistics and Philosophy}, 28(6), 643--686.
\bibitem[Neale(1990)]{Neale1990} Neale, S. (1990). \emph{Descriptions}. MIT Press.
\bibitem[Recanati(2004)]{Recanati2004} Recanati, F. (2004). \emph{Literal Meaning}. Cambridge University Press.
\bibitem[Russell(1905)]{Russell1905} Russell, B. (1905). On denoting. \emph{Mind}, 14(56), 479--493.
\bibitem[Searle(1980)]{Searle1980} Searle, J.~R. (1980). The background of meaning. In J.~R. Searle, F. Kiefer and M. Bierwisch (eds), \emph{Speech Act Theory and Pragmatics}. Reidel.
\bibitem[Stanley(2000)]{Stanley2000} Stanley, J. (2000). Context and logical form. \emph{Linguistics and Philosophy}, 23(4), 391--434.
\bibitem[Stanley(2004)]{Stanley2004} Stanley, J. (2004). On the linguistic basis for contextualism. \emph{Philosophical Studies}, 119(1--2), 119--146.
\bibitem[Stanley and Szab\'o(2000)]{StanleySzabo2000} Stanley, J. and Szab\'o, Z.~G. (2000). On quantifier domain restriction. \emph{Mind \& Language}, 15(2--3), 219--261.
\bibitem[Strawson(1950)]{Strawson1950} Strawson, P.~F. (1950). On referring. \emph{Mind}, 59(235), 320--344.
\bibitem[Travis(1997)]{Travis1997} Travis, C. (1997). Pragmatics. In B. Hale and C. Wright (eds), \emph{A Companion to the Philosophy of Language}. Blackwell.
\bibitem[Wettstein(1981)]{Wettstein1981} Wettstein, H.~K. (1981). Demonstrative reference and definite descriptions. \emph{Philosophical Studies}, 40(2), 241--257.
\end{thebibliography}

\end{document}
